% ==========================================
% SEZIONI DEL VERBALE
% ==========================================
\newpage
\vspace{0.8cm}

\section{Ordine del giorno}

\subsection{Definizione canali comunicazione}
Il gruppo ha scelto di utilizzare per le comunicazioni interne un canale Discord per le riunioni a distanza e Whatsapp per la messaggistica interna.

\subsection{Scelta nome e logo del gruppo}
I membri del gruppo hanno valutato diversi nomi e loghi. A seguito di una votazione, è stato creata la e-mail ufficiale del gruppo.

\subsection{Creazione repository github e Jira come software per il project management}
É stato scelto di creare una repository su GitHub per il versionamento dei file relativi al progetto, inoltre é stato selezionato Jira come software di ausilio all'organizzazione del gruppo. Per la documentazione, invece, è stata creata una cartella condivisa su Google Drive, al fine di garantire a tutti i membri l'accesso, la possibilitá di modifica e la revisione in modalitá sia sincrona che asincrona.

\subsection{Metodologia stesura documenti}
L'approccio scelto per la stesura e la redazione dei documenti é il seguente: dopo una prima bozza redatta su Google Documenti il testo viene approvato dal verificatore. Infine, l'amministratore formatta il documento in LaTex e lo carica sulla repository GitHub.

\subsection{Analisi e discussione capitolati proposti}
Dopo un sondaggio interno, il gruppo ha mostrato un maggiore interesse per i seguenti capitolati:
\begin{enumerate}
    \item \textbf{C1-AURA}: AI for Understanding, Reasoning and Automation
    \item \textbf{C0-AIRA}: AI for Requests Automation
    \item \textbf{C5-RareCare}: Piattaforma AI per la raccolta dati e il supporto decisionale nelle malattie rare
    \item \textbf{C8-Sonar}: Documentazione vivente di uno store Shopify
\end{enumerate}

\subsection{Definizione ruoli}
È stata avviata una discussione preliminare sulla rotazione dei ruoli, che verrà formalizzata in dettaglio nel prossimo incontro.

\subsection{Invio primo contatto con le aziende}
Durante l'incontro odierno abbiamo compilato il form reso disponibile da Vimar per poter partecipare al loro incontro conoscitivo. Inoltre, sono state inviate mail di interesse alle aziende proponenti i capitolati C1, C0 e C5.

\subsection{Programmazione meeting successivo 20 ottobre 2026 15:00}
Il prossimo incontro è stato fissato in data 20/10/2026 sul canale Discord alle ore 15:00 Temi da trattare: rifinire meglio l'assegnazione dei ruoli, discussione approfondita sui capitolati (in particolare sul C1, in preparazione al possibile incontro con l'azienda), visione di gruppo della prima stesura dei documenti assegnati nelle attivitá da svolgere.
